\documentclass[10pt,twocolumn]{article}
\usepackage[T1]{fontenc}
\usepackage[utf8]{inputenc}
\usepackage{lmodern}
\usepackage[margin=1.8cm,top=2.4cm,bottom=2cm,columnsep=0.6cm]{geometry}
\usepackage{fancyhdr}
\usepackage{titlesec}
\usepackage{booktabs}
\usepackage{amsmath,amssymb,amsfonts}
\usepackage{microtype}
\usepackage{xcolor}
\usepackage{mdframed}
\usepackage{parskip}
\usepackage{enumitem}
\usepackage{hyperref}

\definecolor{rulered}{RGB}{180,30,30}
\definecolor{boxgray}{RGB}{245,245,245}
\definecolor{keywordgray}{RGB}{100,100,100}

\pagestyle{fancy}
\fancyhf{}
\fancyhead[L]{\textsc{\small Claude Mag}}
\fancyhead[C]{\small\textit{Machine Learning Research Digest}}
\fancyhead[R]{\small 2026-08-12}
\fancyfoot[C]{\small\thepage}
\renewcommand{\headrulewidth}{0.8pt}

\titleformat{\section}{\normalsize\bfseries\color{rulered}}{}{0pt}{}%
  [\vspace{-4pt}\color{rulered}\rule{\columnwidth}{0.4pt}]
\titlespacing{\section}{0pt}{8pt}{4pt}

\hypersetup{colorlinks=true,urlcolor=rulered,linkcolor=black}

\begin{document}
\setlength{\parindent}{0pt}
\setlength{\parskip}{4pt}

{\small\color{keywordgray}\textbf{Keywords:}
  legged robotics \textbullet{}
  reinforcement learning \textbullet{}
  spatiotemporal control \textbullet{}
  dynamic interception}\par\vspace{4pt}

{\LARGE\bfseries\raggedright
  Spatiotemporal Agility: Time-Constrained Reinforcement
  Learning for Vision-Guided Dynamic Quadrupedal
  Interception}\par\vspace{4pt}

{\large\itshape\raggedright
  Zhejiang University Team Replaces Velocity-Tracking Policies
  with Direct Temporal-Spatial Conditioning to Enable Quadruped
  Ball-Catching Under Strict Deadlines}\par\vspace{6pt}

{\small\textbf{By} Yidong Zhu, Zibo Dai, Tongning Zhang, Leixin Chang, Hua Chen
  (Zhejiang University; LimX Dynamics)}\par\vspace{2pt}

{\small\color{keywordgray}arXiv:2608.06907 \textbullet{} Preprint, August 2026}
\par\vspace{2pt}\noindent{\color{rulered}\rule{\columnwidth}{0.6pt}}\vspace{6pt}

A quadruped robot catching a thrown ball must predict where and when the ball lands
from visual observations, then reach that position before the deadline. A new paper
from Zhejiang University and LimX Dynamics argues that the standard velocity-tracking
locomotion paradigm is fundamentally mismatched to this class of tasks, and proposes
a unified framework that conditions an RL policy directly on spatial targets and
temporal deadlines rather than velocity commands.

\begin{mdframed}[backgroundcolor=boxgray,linecolor=rulered,linewidth=1pt,
    innerleftmargin=6pt,innerrightmargin=6pt,
    innertopmargin=6pt,innerbottommargin=6pt]
\textbf{\small AT A GLANCE}\par\vspace{4pt}
\begin{description}[leftmargin=0pt,labelsep=4pt,itemsep=2pt,parsep=0pt,
    font=\small\bfseries]
  \item[Problem:] Velocity-tracking quadruped policies cannot meet strict temporal
    deadlines for reaching precise positions in dynamic interaction tasks.
  \item[Why it matters:] Dynamic interaction (catching, intercepting, dodging) is the
    next frontier for legged robot agility; no existing policy architecture handles it.
  \item[Method:] Vision module predicts landing point and time; RL policy is conditioned
    on (position, deadline) rather than velocity commands.
  \item[Result:] Ball-catching benchmark significantly outperforms velocity-tracking
    decomposition baselines.
  \item[Evidence:] Evaluated on a standardized ball-catching task with randomized
    trajectory and timing.
\end{description}
\end{mdframed}

\section{The Big Picture}

Legged robots have progressed from laboratory curiosities to practical platforms
over the past decade, largely driven by RL-based locomotion policies that excel at
navigating diverse terrain. But the benchmark tasks that drove this progress---walking
on uneven ground, stair climbing, recovering from pushes---share a common structure:
they reward stable navigation with no hard temporal deadline.

Dynamic interaction tasks are qualitatively different. A robot catching a ball,
intercepting a rolling object, or dodging a thrown projectile must solve a joint
\textbf{spatiotemporal problem}: arrive at a specific location by a specific time.
Velocity-tracking policies, which accept a desired velocity vector and try to track
it, have no native representation of temporal deadlines. Even if they could in
principle navigate to any position, they cannot guarantee arrival by a deadline,
making them unsuitable for dynamic interception.

\section{Why Existing Approaches Fall Short}

The standard decomposition for legged robot tasks is:
\begin{itemize}[itemsep=2pt,parsep=0pt]
  \item \textbf{High-level planner:} computes a path and velocity profile
  \item \textbf{Low-level controller:} tracks the desired velocity command
\end{itemize}

This decomposition fails for time-constrained interception for two reasons.
First, velocity-tracking introduces an abstraction layer that discards the
temporal deadline: the planner must guess a velocity profile that achieves
arrival by the deadline, but the tracking error of the low-level controller
means the actual arrival time is unpredictable.

Second, the two-level decomposition requires the planner to predict the
controller's behavior, introducing compounding errors. In practice, systems
that work well for slow, navigational tasks become brittle when deadlines
are tight and the target is moving.

\section{Core Mechanism}

\textbf{Vision Module.}
A neural vision module processes monocular camera images to predict:
\begin{itemize}[itemsep=2pt,parsep=0pt]
  \item $\hat{\mathbf{p}} \in \mathbb{R}^3$: predicted landing position (world frame)
  \item $\hat{\tau} \in \mathbb{R}_{>0}$: predicted time-to-landing (seconds)
\end{itemize}
The module is trained on ball-trajectory data with supervised regression.

\textbf{Time-Conditioned RL Policy.}
Rather than consuming velocity commands, the locomotion policy is conditioned on the
pair $(\hat{\mathbf{p}}, \hat{\tau})$. This directly expresses the spatiotemporal
constraint: reach position $\hat{\mathbf{p}}$ within $\hat{\tau}$ seconds. The
policy must learn internally to compute appropriate gaits, step frequencies, and
postures to meet both the spatial and temporal objectives simultaneously.

\section{Technical Formulation}

Let the robot state at time $t$ be $s_t$ and the action (joint torques or velocity
targets) be $a_t$. The policy $\pi_\theta(a_t \mid s_t, \hat{\mathbf{p}}, \hat{\tau} - t)$
is conditioned on the remaining time budget $\hat{\tau} - t$, which decreases with
each timestep, providing an implicit urgency signal.

The reward is:
\[
r_t = \lambda_{\mathrm{pos}}\,r_{\mathrm{pos}}(s_T, \hat{\mathbf{p}})
    + \lambda_{\mathrm{time}}\,r_{\mathrm{time}}(T, \hat{\tau})
    + \lambda_{\mathrm{stab}}\,r_{\mathrm{stab}}(s_t)
\]
where $T$ is the actual interception time, $r_{\mathrm{pos}}$ rewards proximity
to the landing point, $r_{\mathrm{time}}$ penalizes arriving after the deadline,
and $r_{\mathrm{stab}}$ encourages stable locomotion.

The vision module and locomotion policy are trained jointly, with the RL reward
backpropagating a signal to align vision predictions with achievable interception
targets.

\section{Learning or Inference Procedure}

\textbf{Simulation training:} The robot and ball physics are simulated with
randomized launch parameters (speed, angle, direction) to cover diverse
trajectory families. Domain randomization is applied to robot dynamics to
improve sim-to-real transfer.

\textbf{Curriculum:} Training begins with slow, predictable ball trajectories
and gradually increases speed and variability as the policy improves.

\textbf{Inference:} The vision module processes each incoming frame, updating
$\hat{\mathbf{p}}$ and $\hat{\tau}$ in real time. The locomotion policy uses
the latest prediction to compute joint targets at 50~Hz.

\section{What the Guarantee Says}

No formal guarantee is provided. The policy's behavior near the deadline depends
on learned heuristics rather than provable time-optimal control. However, the
experimental results demonstrate that the time-conditioned formulation produces
significantly more reliable deadline adherence than velocity-tracking baselines,
suggesting the implicit urgency signal is effective in practice.

\section{Experimental Findings}

\begin{itemize}[itemsep=2pt,parsep=0pt]
  \item \textbf{Ball-catching success rate:} The time-conditioned RL policy
    significantly outperforms velocity-tracking decomposition baselines on
    the proposed benchmark.
  \item \textbf{Deadline adherence:} The policy learns to vary gait speed
    continuously based on remaining time, rather than running at maximum
    speed and stopping.
  \item \textbf{Trajectory diversity:} Performance degrades gracefully with
    increasing ball speed and trajectory variability.
  \item \textbf{Sim-to-real:} The framework transfers to physical hardware
    with domain randomization, demonstrating real-world viability.
\end{itemize}

\section{Ablations and Interpretation}

\textbf{Time conditioning vs.\ position-only conditioning.}
Removing the time dimension and conditioning only on target position degrades
success rate substantially, confirming that the temporal deadline is a necessary
conditioning input rather than information the policy can infer from position alone.

\textbf{Joint vs.\ separate training.}
Training the vision module and locomotion policy separately (first vision, then RL
with fixed vision) produces lower success rates than joint training, suggesting
that the locomotion policy's behavior affects what landing predictions are achievable.

\textbf{Curriculum vs.\ full distribution.}
Training without curriculum (full ball-speed distribution from the start) results
in slower convergence and lower asymptotic performance, particularly for fast
trajectories.

\section*{Reference}

Yidong Zhu, Zibo Dai, Tongning Zhang, Leixin Chang, Hua Chen.
\textbf{Spatiotemporal Agility: Time-Constrained Reinforcement Learning for
Vision-Guided Dynamic Quadrupedal Interception}.
arXiv:2608.06907, August 2026.
\url{https://arxiv.org/abs/2608.06907}

\end{document}
